% AgriFlow — Audit Menyeluruh Repo, GitHub, dan Deployment
% Compile: pdflatex AgriFlow_Audit_2026-07.tex  (jalankan 2x untuk ToC)
\documentclass[11pt,a4paper]{article}

\usepackage[utf8]{inputenc}
\usepackage[T1]{fontenc}
\usepackage[a4paper,margin=2.3cm]{geometry}
\usepackage{amsmath,amssymb}
% NOTE: jangan muat longtable maupun xcolor[table]/colortbl di sini. Pada
% instalasi MiKTeX ini colortbl merusak template kolom p{} (\insert@pcolumn
% undefined), sehingga zebra-striping sengaja tidak dipakai. Cukup booktabs.
\usepackage{booktabs}
\usepackage{array}
\usepackage{enumitem}
\usepackage{xcolor}
\usepackage[hidelinks]{hyperref}
\usepackage{titlesec}
\usepackage{fancyhdr}
\usepackage{listings}

\definecolor{agrigreen}{HTML}{1B5E20}
\definecolor{sevhigh}{HTML}{C62828}
\definecolor{sevmed}{HTML}{EF6C00}
\definecolor{sevlow}{HTML}{1565C0}
\definecolor{okgreen}{HTML}{2E7D32}

\titleformat{\section}{\Large\bfseries\color{agrigreen}}{\thesection}{0.6em}{}
\titleformat{\subsection}{\large\bfseries\color{black!80}}{\thesubsection}{0.6em}{}
\setlist{nosep,leftmargin=1.3em}

\newcommand{\code}[1]{\texttt{\small #1}}
\newcommand{\sevH}{\textcolor{sevhigh}{\textbf{TINGGI}}}
\newcommand{\sevM}{\textcolor{sevmed}{\textbf{SEDANG}}}
\newcommand{\sevL}{\textcolor{sevlow}{\textbf{RENDAH}}}
\newcommand{\ok}{\textcolor{okgreen}{\textbf{LULUS}}}
\newcommand{\gagal}{\textcolor{sevhigh}{\textbf{GAGAL}}}

\lstset{
  basicstyle=\ttfamily\footnotesize,
  breaklines=true,
  frame=single,
  rulecolor=\color{black!25},
  backgroundcolor=\color{black!3},
  columns=fullflexible,
}

% Identifier panjang ber-underscore (mis. get_active_demand_event) memicu
% overfull hbox. Longgarkan justifikasi dan beri ruang darurat.
\sloppy
\emergencystretch=3em

\pagestyle{fancy}
\fancyhf{}
\fancyhead[L]{\small AgriFlow — Audit Menyeluruh}
\fancyhead[R]{\small Juli 2026}
\fancyfoot[C]{\small \thepage}
\renewcommand{\headrulewidth}{0.3pt}

\title{\textbf{AgriFlow — Audit Menyeluruh}\\[4pt]
\large Fitur, Repositori GitHub, Deployment Web, dan Rekonsiliasi Audit Sebelumnya}
\author{Hilmi \\ \small\url{https://master-hilmi.vercel.app/}}
\date{17 Juli 2026}

\begin{document}
\maketitle

\begin{abstract}
\noindent
Dokumen ini adalah audit menyeluruh terhadap proyek AgriFlow per 17 Juli 2026, mencakup tiga
permukaan sekaligus: (1) kode dan fitur di repositori lokal, (2) repositori GitHub beserta
continuous integration, dan (3) deployment web yang berjalan (dashboard Vercel dan API Hugging
Face Space). Setiap klaim di dokumen ini diverifikasi dengan menjalankan kode, memanggil endpoint
produksi, atau membaca log CI, bukan dengan membaca dokumentasi yang ada. Temuan utama: seluruh
404 tes lulus di mesin lokal, namun CI di \code{main} berstatus merah pada keempat matrix leg
karena dua tes bergantung pada tanggal kalender sistem. Akar masalahnya bukan sekadar tes rapuh,
melainkan cacat logika di mesin pencocokan yang menyebabkan event kalender membatalkan flag
kebijakan yang diberikan pemanggil secara eksplisit. Audit juga menemukan dua ketidakcocokan
antara klaim dokumentasi dan artefak yang benar-benar dikirim, serta memastikan tidak ada
kebocoran kredensial. Dokumen \code{AUDIT\_v10.md} yang ada dinyatakan usang dan digantikan oleh
dokumen ini.
\end{abstract}

\tableofcontents
\vspace{1em}
\hrule
\vspace{1em}

% =====================================================================
\section{Ringkasan Eksekutif}

\subsection{Verdict}

Inti teknis AgriFlow (mesin pencocokan 4 lapis, deteksi anomali, penyajian via API dan dashboard)
\textbf{berfungsi dan terbukti berjalan}. Rantai dashboard sampai API tervalidasi hidup
end-to-end. Yang bermasalah adalah \textbf{lapisan integritas}: CI merah, dan dua klaim di
dokumentasi lebih besar daripada artefak yang benar-benar dikirim.

\begin{center}
\begin{tabular}{@{}p{5.4cm}cp{7.6cm}@{}}
\toprule
\textbf{Permukaan audit} & \textbf{Status} & \textbf{Catatan} \\
\midrule
Test suite (lokal) & \ok & 403 lulus, 1 di-skip, 41,9 detik \\
Test suite (CI GitHub) & \gagal & 2 gagal $\times$ 4 matrix leg, bergantung tanggal \\
Mesin pencocokan 4 lapis & \ok & Berjalan di data BPS asli 2022, 63 match \\
Deteksi anomali harga & \ok & Berjalan, tetapi salah label metode \\
Prediksi harga & \ok & Berjalan sebagai baseline, bukan TimesFM \\
API produksi (HF Space) & \ok & 6 endpoint hidup, CORS benar \\
Dashboard produksi (Vercel) & \ok & HTTP 200, terhubung ke API produksi \\
Latency benchmark & \ok & p99 tertinggi 59,53 ms, margin 88,1\,\% \\
Kebersihan kredensial & \ok & Tidak ada secret ter-commit, riwayat bersih \\
Konsistensi klaim vs bukti & \gagal & 2 overclaim di README (lihat \S\ref{sec:klaim}) \\
\bottomrule
\end{tabular}
\end{center}

\subsection{Tiga hal yang paling perlu ditindaklanjuti}

\begin{enumerate}
  \item \textbf{CI merah di \code{main} sejak 14 Juli 2026} akibat cacat logika kalender di
        \code{engine.py}, bukan sekadar tes rapuh. Ini akan gagal ulang otomatis setiap
        1--15 Januari dan 1--15 Juli (sekitar 30 hari per tahun). Lihat \S\ref{sec:f1}.
  \item \textbf{README mengklaim forecasting memakai TimesFM 2.0}, padahal 56 dari 56 ramalan
        yang dikirim berlabel \code{seasonal\_naive\_baseline}. API produksi justru jujur
        melaporkan label baseline. Lihat \S\ref{sec:f2}.
  \item \textbf{Detektor anomali menamai dirinya S-H-ESD} padahal implementasinya Hampel/MAD
        tanpa uji ESD sama sekali. API produksi masih menyiarkan \code{"method":"shesd\_v2"}.
        Perbaikannya sudah ada tetapi terhenti di branch lokal yang belum di-push.
        Lihat \S\ref{sec:f3}.
\end{enumerate}

% =====================================================================
\section{Metodologi dan Cakupan Audit}

Audit ini memverifikasi dengan eksekusi, bukan dengan pembacaan dokumen. Yang benar-benar
dijalankan atau dipanggil:

\begin{center}
\begin{tabular}{@{}p{6.8cm}p{8.2cm}@{}}
\toprule
\textbf{Tindakan verifikasi} & \textbf{Sumber bukti} \\
\midrule
Menjalankan seluruh test suite & \code{python -m pytest tests/ -q} \\
Menjalankan demo data nyata & \code{python examples/run\_demo\_real.py} \\
Menjalankan benchmark latency & \code{python benchmarks/latency.py} \\
Membaca log CI yang gagal & \code{gh run view 29348620529 -{}-log-failed} \\
Reproduksi bug kalender & Pemanggilan langsung \code{get\_active\_demand\_event()} \\
Memanggil 6 endpoint API produksi & \code{curl} ke Hugging Face Space \\
Memeriksa wiring dashboard & Grep bundel JS hasil build di Vercel \\
Memeriksa CORS & Preflight \code{OPTIONS} dari origin Vercel \\
Memeriksa kebocoran secret & \code{git ls-files} dan pemindaian riwayat commit \\
\bottomrule
\end{tabular}
\end{center}

\noindent
\textbf{Lingkungan.} Branch \code{main} pada commit \code{a6b3b47}, Python 3.14.3 (venv lokal),
Windows 11. CI GitHub Actions memakai matriks Ubuntu dan Windows $\times$ Python 3.11 dan 3.12.

% =====================================================================
\section{Inventaris Fitur yang Terverifikasi}

Repositori berisi 6.743 baris kode Python produksi (di luar tes) yang terbagi ke lima pilar.

\subsection{Pilar 1: Distribusi (mesin pencocokan 4 lapis)}

Modul \code{matching\_engine/} berisi 2.101 baris di lima file.

\begin{center}
\begin{tabular}{@{}p{3.6cm}rp{9.6cm}@{}}
\toprule
\textbf{Modul} & \textbf{Baris} & \textbf{Peran} \\
\midrule
\code{engine.py}      & 653 & Orkestrator, kalender event, Bulog reserve, blackout rute \\
\code{allocation.py}  & 467 & Lapis 3, alokasi dan pengali equity berbasis IPM \\
\code{constraints.py} & 453 & Lapis 1, 8 hard constraint dan pembangkit kandidat \\
\code{scoring.py}     & 270 & Lapis 2, skor 5 dimensi dan 6 profil bobot \\
\code{models.py}      & 258 & Tipe domain (SupplyNode, DemandNode, Match) \\
\bottomrule
\end{tabular}
\end{center}

\noindent
\textbf{Bukti eksekusi} (\code{run\_demo\_real.py}, data BPS asli 2022):

\begin{center}
\begin{tabular}{@{}lrr@{}}
\toprule
\textbf{Komoditas} & \textbf{Match} & \textbf{Volume tercocokkan (ton)} \\
\midrule
beras\_premium & 12 & 222.386,6 \\
beras\_medium  & 12 & 148.257,8 \\
cabai\_merah   & 21 & 22.737,5 \\
cabai\_rawit   & 18 & 21.828,7 \\
bawang\_putih  & 0  & Tidak ada surplus domestik, dirutekan ke saran impor \\
\midrule
\textbf{Total} & \textbf{63} & \textbf{415.210,6} \\
\bottomrule
\end{tabular}
\end{center}

\noindent
Kasus \code{bawang\_putih} patut dicatat sebagai perilaku benar, bukan kegagalan: produksi Jatim
2022 sekitar 855 ton berbanding konsumsi sekitar 80.640 ton, sehingga seluruh 38 kabupaten
berstatus defisit. Mesin menanganinya dengan anggun (0 match, lalu mengeluarkan
\code{external\_opportunity} berupa saran impor).

\subsection{Pilar 2: Deteksi anomali harga}

\code{analysis/price\_anomaly.py} (513 baris) melakukan deseasonalisasi lalu menerapkan MAD
bergulir pada residual, dengan MAD floor untuk mencegah pemicuan trivial saat dispersi sangat
kecil. Terverifikasi berjalan di produksi (endpoint mengembalikan 50 anomali). Catatan penting
soal penamaan metode ada di \S\ref{sec:f3}.

\subsection{Pilar 3: Prediksi harga}

\code{analysis/forecast\_timesfm.py} (309 baris) memprecompute ramalan secara offline; server
tidak pernah meng-import TimesFM saat runtime (penjaga OOM di HF Space free tier). Artefak yang
dikirim, \code{sample\_data/forecasts/forecast\_all.json}:

\begin{center}
\begin{tabular}{@{}p{4.6cm}p{10.4cm}@{}}
\toprule
\textbf{Properti} & \textbf{Nilai terukur} \\
\midrule
Jumlah rekaman & 56 (7 komoditas $\times$ 8 kota) \\
Label metode & \code{seasonal\_naive\_baseline} pada 56 dari 56 rekaman \\
Horizon & 30 hari \\
Akhir riwayat & 2025-12-31 \\
Interval & p10 dan p90 per titik \\
\bottomrule
\end{tabular}
\end{center}

\subsection{Pilar 4: Akses (API, dashboard, bot WhatsApp)}

\code{whatsapp\_bot/server.py} (495 baris) melayani webhook WhatsApp sekaligus enam endpoint REST
yang dikonsumsi dashboard. Dashboard adalah aplikasi Next.js di \code{dashboard/}. Bot memiliki
pemisahan intent (\code{intent.py}), handler (\code{handlers.py}), klien Gemini
(\code{gemini\_client.py}), dan klien Twilio (\code{twilio\_client.py}).

\subsection{Pilar 5: Konektor data dan persistensi}

Delapan konektor di \code{data\_sources/} (1.296 baris), seluruhnya dual-mode (mock dan live):
\code{bps.py}, \code{pihps\_bi.py}, \code{bapanas.py}, \code{bmkg.py}, \code{bnpb.py},
\code{pvmbg.py}, \code{google\_maps.py}, \code{hijri\_calendar.py}. Persistensi di \code{db/}
(458 baris) berisi \code{schema.sql}, \code{db\_loader.py}, dan \code{price\_ingest.py}.

% =====================================================================
\section{Audit Test Suite}

\subsection{Hasil lokal}

\begin{lstlisting}
$ python -m pytest tests/ -q
403 passed, 1 skipped in 41.94s
\end{lstlisting}

\noindent
Sebaran 404 tes yang terkumpul di 21 file:

\begin{center}
\begin{tabular}{@{}lr@{\hskip 1.2cm}lr@{}}
\toprule
\textbf{File} & \textbf{Tes} & \textbf{File} & \textbf{Tes} \\
\midrule
test\_price\_anomaly            & 49 & test\_layer0\_tier                  & 16 \\
test\_forecast\_anomaly\_api     & 40 & test\_constrained\_scenario         & 15 \\
test\_price\_ingest             & 33 & test\_layer3\_allocation            & 14 \\
test\_scenarios\_tier1\_ext      & 27 & test\_scenarios\_political          & 11 \\
test\_layer2\_scoring           & 24 & test\_db\_loader                    & 11 \\
test\_baseline\_comparison      & 24 & test\_scenarios\_disruption         & 9 \\
test\_real\_data\_pipeline       & 23 & test\_road\_distance                & 9 \\
test\_real\_sd\_horti\_2022       & 22 & test\_scenarios\_temporal           & 7 \\
test\_real\_sd\_beras            & 21 & test\_scenarios\_spatial            & 6 \\
test\_whatsapp\_bot             & 20 & test\_scenarios\_volume             & 4 \\
test\_layer1\_constraints       & 19 & & \\
\bottomrule
\end{tabular}
\end{center}

\subsection{Hasil CI: divergensi yang menyingkap bug}

CI pada commit \code{a6b3b47} gagal di \textbf{keempat} matrix leg dengan pola identik:

\begin{lstlisting}
2 failed, 394 passed, 8 skipped, 1 warning in 34.90s
FAILED tests/test_scenarios_political.py::TestE4_ImportPolicy::test_import_policy_uses_alternate_weights
FAILED tests/test_scenarios_temporal.py::TestC1_RamadanSpike::test_ramadan_active_uses_ramadan_weights
\end{lstlisting}

\noindent
Divergensi lokal (lulus) versus CI (gagal) pada commit yang sama inilah petunjuk yang membawa ke
temuan F1. Penjelasan lengkapnya di \S\ref{sec:f1}.

% =====================================================================
\section{Audit Repositori GitHub}

\begin{center}
\begin{tabular}{@{}p{4.2cm}p{10.8cm}@{}}
\toprule
\textbf{Atribut} & \textbf{Nilai} \\
\midrule
URL & \url{https://github.com/masterA88/agriflow-engine} \\
Visibilitas & Publik \\
Push terakhir & 2026-07-14T16:11:41Z \\
Commit \code{main} & \code{a6b3b47} (Merge PR \#1: update dashboard) \\
Pull request & 1 total, berstatus merged \\
Issue terbuka & 0 \\
Lisensi & MIT \\
\bottomrule
\end{tabular}
\end{center}

\subsection{Continuous integration}

Workflow \code{.github/workflows/test.yml} menjalankan pytest pada matriks
\{Ubuntu, Windows\} $\times$ \{Python 3.11, 3.12\}, ditambah job benchmark latency yang
mempublikasikan artefak tanpa menjadi gate.

\begin{center}
\begin{tabular}{@{}p{7.4cm}p{2.4cm}p{4.4cm}@{}}
\toprule
\textbf{Commit} & \textbf{Status} & \textbf{Tanggal} \\
\midrule
Merge PR \#1: update dashboard & \gagal & 2026-07-14 \\
docs: embed research poster in README & \ok & 2026-06-03 \\
feat(forecast): reproducible holdout backtest & \ok & 2026-06-03 \\
feat(forecast): vendor TimesFM 2.0 pipeline & \ok & 2026-06-03 \\
docs: add AgriFlow logo to README & \ok & 2026-06-02 \\
\bottomrule
\end{tabular}
\end{center}

\noindent
\textbf{Pengamatan.} CI hijau sepanjang Juni dan baru merah pada 14 Juli. Ini konsisten dengan
akar masalah berbasis tanggal (\S\ref{sec:f1}), bukan dengan regresi yang diperkenalkan PR \#1.
PR \#1 hanya menyentuh dashboard dan tidak menyentuh \code{matching\_engine/} sama sekali.

\subsection{Divergensi branch}

\begin{center}
\begin{tabular}{@{}p{4.2cm}p{10.8cm}@{}}
\toprule
\textbf{Branch} & \textbf{Kondisi} \\
\midrule
\code{main} & Sinkron dengan \code{origin/main} pada \code{a6b3b47} \\
\code{post-submit-analysis} & \textbf{Lokal saja, belum di-push.} 21 file, +2.886/-2.046 baris \\
\code{space-deploy-v3} & Lokal, jalur deploy Hugging Face \\
\code{Chelseaayu/main} & Lokal, kontribusi rekan tim \\
\bottomrule
\end{tabular}
\end{center}

\noindent
Branch \code{post-submit-analysis} memuat kerja bernilai yang belum terlindungi remote, termasuk
\code{analysis/backtest\_baseline.py}, kalibrasi interval, koreksi pelabelan detektor anomali,
harness sensitivitas bobot, dan pipeline TimesFM di \code{prediction/timesfm/}.

% =====================================================================
\section{Audit Deployment Web}

\subsection{Dashboard (Vercel)}

\begin{center}
\begin{tabular}{@{}p{4.6cm}p{10.4cm}@{}}
\toprule
\textbf{Pemeriksaan} & \textbf{Hasil} \\
\midrule
URL & \url{https://agriflow-engine.vercel.app/} \\
Status HTTP & 200, waktu muat 1,18 detik \\
Judul dokumen & \code{AgriFlow Dashboard} \\
Basis API pada bundel build & \code{https://masterAAA123-agriflow-api.hf.space} \\
Fallback \code{localhost:8000} & Ada di bundel, tetapi tidak aktif karena env ter-set \\
\bottomrule
\end{tabular}
\end{center}

\noindent
\code{dashboard/app/lib/api.ts} memakai pola \code{process.env.NEXT\_PUBLIC\_API\_URL ||
"http://localhost:8000"}. Pola ini berisiko diam-diam gagal bila variabel tidak ter-set saat
build. Audit memverifikasi bahwa variabel \textbf{ter-set dengan benar} di Vercel, sehingga
risiko tersebut tidak terwujud. Ini dicatat sebagai kerapuhan laten, bukan cacat aktif.

\subsection{API (Hugging Face Space)}

\begin{center}
\begin{tabular}{@{}p{6.4cm}cp{7.0cm}@{}}
\toprule
\textbf{Endpoint} & \textbf{Kode} & \textbf{Respons terverifikasi} \\
\midrule
\code{/health}                  & 200 & \code{status:ok}, 38 kabupaten, 19 komoditas \\
\code{/api/v1/commodities}      & 200 & Daftar komoditas \\
\code{/api/v1/kabupaten}        & 200 & 38 kabupaten dengan lat, lng, tier, IPM \\
\code{/api/v1/surplus-deficit}  & 200 & Baris surplus dan defisit plus total \\
\code{/api/v1/matches}          & 200 & 23 match, terskor dan terurut \\
\code{/api/v1/forecast}         & 200 & Titik, p10, p90 untuk 30 hari \\
\code{/api/v1/anomalies}        & 200 & 50 anomali \\
\code{/} (root)                 & 404 & Wajar, tidak ada route root \\
\bottomrule
\end{tabular}
\end{center}

\noindent
\textbf{CORS terverifikasi benar.} Preflight \code{OPTIONS} dari origin Vercel mengembalikan
\code{access-control-allow-origin: https://agriflow-engine.vercel.app} dengan
\code{access-control-allow-methods: GET}. Origin di-whitelist secara spesifik, bukan wildcard.

\noindent
\textbf{Rantai end-to-end tervalidasi hidup:} browser memuat dashboard Vercel, bundel menunjuk ke
host HF Space yang benar, host merespons 200, dan CORS mengizinkan origin tersebut.

% =====================================================================
\section{Temuan}

\begin{center}
\begin{tabular}{@{}clp{8.6cm}c@{}}
\toprule
\textbf{ID} & \textbf{Severity} & \textbf{Ringkasan} & \textbf{Bagian} \\
\midrule
F1 & \sevH & Event kalender membatalkan flag kebijakan eksplisit, CI merah & \ref{sec:f1} \\
F2 & \sevH & README klaim TimesFM, artefak 100\,\% baseline & \ref{sec:f2} \\
F3 & \sevM & Detektor anomali salah label sebagai S-H-ESD & \ref{sec:f3} \\
F4 & \sevM & Kerja bernilai terjebak di branch lokal yang belum di-push & \ref{sec:f4} \\
F5 & \sevM & \code{AUDIT\_v10.md} usang dan menyesatkan & \ref{sec:f5} \\
F6 & \sevL & API produksi berjalan dalam \code{mock\_mode} & \ref{sec:f6} \\
F7 & \sevL & Parameter \code{city} pada endpoint forecast tidak intuitif & \ref{sec:f7} \\
\bottomrule
\end{tabular}
\end{center}

% ---------------------------------------------------------------
\subsection{F1 (\sevH) Event kalender membatalkan flag kebijakan eksplisit}
\label{sec:f1}

\subsubsection*{Gejala}

Dua tes gagal di keempat matrix leg CI pada 14 Juli 2026, tetapi lulus di mesin lokal pada 17 Juli
2026 dengan commit yang sama.

\subsubsection*{Akar masalah}

\code{matching\_engine/engine.py:246-248} mendefinisikan window \code{SCHOOL\_START} sebagai 14
hari sebelum 15 Juli dan 15 Januari, yaitu 1--15 Juli dan 1--15 Januari. CI berjalan 14 Juli,
tepat di dalam window. Audit mereproduksinya secara langsung:

\begin{lstlisting}
>>> get_active_demand_event(datetime(2026, 7, 14))   # tanggal CI
'SCHOOL_START'
>>> get_active_demand_event(datetime(2026, 7, 17))   # tanggal audit lokal
None
\end{lstlisting}

\noindent
Kode pemilihan bobot di \code{engine.py:449-468}:

\begin{lstlisting}
active_event = get_active_demand_event(reference_date)
if logistics.is_ramadan_proximity:
    active_event = active_event or "RAMADAN"   # BUG: "or" kalah dari SCHOOL_START
ramadan_active = active_event == "RAMADAN"
if active_event:
    weights = _weights_for_event(active_event)
elif import_policy_active:                     # BUG: tak pernah tercapai saat event aktif
    weights = scoring.IMPORT_POLICY_WEIGHTS
else:
    weights = scoring.DEFAULT_WEIGHTS
\end{lstlisting}

\noindent
Bukti konfirmasi paling kuat: pesan error CI melaporkan bobot aktual berisi \code{price: 0.22},
dan \code{SCHOOL\_START\_WEIGHTS} memang bernilai \code{price: 0.22}, sedangkan
\code{IMPORT\_POLICY\_WEIGHTS} yang diharapkan bernilai \code{price: 0.10}.

\subsubsection*{Mengapa ini bug produksi, bukan sekadar tes rapuh}

\begin{enumerate}
  \item \textbf{Implementasi melanggar prioritas yang didokumentasikannya sendiri.} Docstring
        \code{get\_active\_demand\_event} menyatakan ``Priority order saat overlap: RAMADAN >
        NATAL > IMLEK > SCHOOL\_START''. Namun ketika pemanggil menegaskan Ramadan lewat
        \code{logistics.is\_ramadan\_proximity}, operator \code{or} hanya mengisi nilai bila
        \code{active\_event} bernilai \code{None}. Akibatnya \code{SCHOOL\_START} yang diturunkan
        dari tanggal justru \textbf{mengalahkan} Ramadan yang dinyatakan eksplisit, persis
        kebalikan dari prioritas yang ditulis.
  \item \textbf{Kebijakan impor menjadi tidak dapat diaktifkan selama 30 hari per tahun.}
        Struktur \code{elif} membuat \code{import\_policy\_active=True} diabaikan diam-diam
        sepanjang window event mana pun. Kebijakan impor adalah kondisi regulatif yang ortogonal
        terhadap event musiman; keduanya seharusnya dapat berlaku bersamaan.
  \item \textbf{Menghasilkan state yang saling bertentangan.} Pada \code{engine.py:573}, flag
        \code{IMPORT\_POLICY\_ACTIVE} \textbf{tetap dipasang} ke setiap match meskipun bobot
        import policy tidak pernah diterapkan. Konsumen hilir (dashboard, audit trail) akan
        membaca match yang mengaku tunduk pada kebijakan impor padahal diskor dengan bobot
        kos-kosan awal tahun ajaran. Untuk sistem B2G yang menjual transparansi kepada Bank
        Indonesia dan Pemda, jejak audit yang berbohong adalah cacat serius.
\end{enumerate}

\subsubsection*{Dampak}

CI akan merah secara otomatis sekitar 30 hari setiap tahun (1--15 Januari dan 1--15 Juli) tanpa
ada perubahan kode. Lebih jauh, setiap keputusan alokasi yang dijalankan dalam window tersebut
memakai profil bobot yang salah bila pemanggil mengaktifkan kebijakan impor atau menegaskan
Ramadan.

\subsubsection*{Perbaikan yang disarankan}

\begin{enumerate}
  \item \textbf{Hormati override eksplisit dari pemanggil}, sesuai prioritas terdokumentasi:
\begin{lstlisting}
active_event = get_active_demand_event(reference_date)
if logistics.is_ramadan_proximity:
    active_event = "RAMADAN"   # penegasan pemanggil menang, sesuai docstring
\end{lstlisting}
  \item \textbf{Pisahkan kebijakan impor dari event kalender.} Kebijakan impor sebaiknya
        menyesuaikan bobot secara komposisional di atas profil event, atau minimal memancarkan
        warning eksplisit ketika ditekan oleh sebuah event, sehingga tidak pernah gagal secara
        diam-diam.
  \item \textbf{Sematkan \code{reference\_date} di tes.} \code{run\_matching()} sudah menerima
        parameter \code{reference\_date}, sehingga kedua tes dapat dibuat deterministik tanpa
        mengubah kode produksi. Ini menutup kerapuhan tes, tetapi \textbf{jangan} dijadikan
        satu-satunya perbaikan: menyematkan tanggal saja akan menyembunyikan bug produksi di
        butir 1 dan 2.
\end{enumerate}

% ---------------------------------------------------------------
\subsection{F2 (\sevH) Klaim TimesFM versus artefak baseline}
\label{sec:f2}

README menyatakan pada tabel fitur berstatus centang hijau:

\begin{quote}
``\textbf{Prediksi} — Forecasting harga 30 hari dengan \textbf{TimesFM 2.0} (foundation model
time-series)''
\end{quote}

\noindent
Verifikasi terhadap artefak yang benar-benar dikirim menunjukkan sebaliknya:

\begin{center}
\begin{tabular}{@{}p{6.4cm}p{8.6cm}@{}}
\toprule
\textbf{Sumber} & \textbf{Label metode terukur} \\
\midrule
\code{forecast\_all.json} (ter-commit) & \code{seasonal\_naive\_baseline}, 56 dari 56 rekaman \\
API produksi (\code{/api/v1/forecast}) & \code{"method":"seasonal\_naive\_baseline"} \\
README & ``TimesFM 2.0'' dengan status centang \\
\bottomrule
\end{tabular}
\end{center}

\noindent
Ironinya, \textbf{kode dan API sudah jujur}; hanya README yang melebihkan. Fallback baseline
diberi label apa adanya karena model 2 GB mengalami OOM di HF Space free tier, dan itu justru
keputusan rekayasa yang benar. Yang perlu diperbaiki adalah narasinya.

\noindent
\textbf{Konteks pendukung.} Backtest holdout 30 hari yang reproducible (di branch
\code{post-submit-analysis}) mengukur MAPE keseluruhan baseline sebesar 10,8\,\%. Klaim TimesFM
MAPE 3,0\,\% yang beredar di dokumen proposal tidak dapat direproduksi dari file mana pun di
repositori ini.

\noindent
\textbf{Perbaikan.} Ubah baris README menjadi label yang akurat, misalnya ``Forecasting harga 30
hari dengan seasonal naive baseline (MAPE 10,8\,\% pada holdout 30 hari); pipeline TimesFM 2.0
tervendor dan menunggu runtime ber-GPU''. Ini tetap kuat, dan tidak dapat dibongkar di sesi tanya
jawab teknis.

% ---------------------------------------------------------------
\subsection{F3 (\sevM) Detektor anomali salah label sebagai S-H-ESD}
\label{sec:f3}

\code{analysis/price\_anomaly.py} pada \code{main} mendeskripsikan dirinya sebagai
``S-H-ESD price anomaly detector'' dan ``METHOD: Seasonal-Hybrid ESD (S-H-ESD style)''. Namun
implementasi aktual melakukan deseasonalisasi lalu menerapkan MAD bergulir pada residual, yaitu
detektor Hampel. \textbf{Tidak ada uji Extreme Studentized Deviate sama sekali} di dalam kode.
API produksi meneruskan label keliru ini ke konsumen: \code{"method":"shesd\_v2"}.

\noindent
Perbaikan pelabelan sudah ditulis, tetapi hanya ada di branch \code{post-submit-analysis} yang
belum di-push, sehingga produksi masih menyiarkan label yang salah.

\noindent
\textbf{Mengapa penting.} Hampel/MAD adalah metode yang sah, robust, dan tepat untuk kasus ini.
Tidak ada alasan teknis untuk menyebutnya S-H-ESD. Salah label terhadap juri akademis justru
menciptakan risiko kredibilitas yang tidak sepadan, padahal metode aslinya sudah dapat dibela.

% ---------------------------------------------------------------
\subsection{F4 (\sevM) Kerja bernilai terjebak di branch lokal}
\label{sec:f4}

Branch \code{post-submit-analysis} menyimpan 2.886 baris tambahan di 21 file, mencakup backtest
reproducible, kalibrasi interval p10/p90 dari 1,4826$\times$MAD menjadi 1,90$\times$MAD, koreksi
pelabelan F3, harness sensitivitas bobot, dan pipeline TimesFM. Seluruhnya hanya ada di satu
mesin, tanpa salinan remote.

\noindent
\textbf{Catatan tata kelola.} Branch ini sengaja tidak di-push, dan audit ini
\textbf{tidak mengubah statusnya}. Push memerlukan otorisasi eksplisit dari pemilik repositori.
Rekomendasinya adalah mempertimbangkan push sebagai branch (tanpa merge) semata-mata untuk
mendapatkan cadangan remote.

% ---------------------------------------------------------------
\subsection{F5 (\sevM) Audit sebelumnya sudah usang}
\label{sec:f5}

\code{docs/AUDIT\_v10.md} (4 Mei 2026) kini menyesatkan bila dibaca sebagai kondisi terkini.

\begin{center}
\begin{tabular}{@{}p{5.0cm}p{4.6cm}p{5.0cm}@{}}
\toprule
\textbf{Aspek} & \textbf{AUDIT\_v10 (Mei)} & \textbf{Terukur (Juli)} \\
\midrule
Jumlah tes & 106 & 404 terkumpul, 403 lulus \\
Runtime tes & 0,16 detik & 41,94 detik \\
File tes & 10 & 21 \\
Cakupan & Mesin pencocokan saja & Ditambah deteksi, prediksi, API, bot \\
Status CI & Belum ada (dicatat sebagai open question) & Ada, sedang merah \\
\bottomrule
\end{tabular}
\end{center}

\noindent
Beberapa isi \code{AUDIT\_v10.md} tetap berharga dan sebaiknya dipertahankan, khususnya analisis
bottleneck skala nasional (\S4), tinjauan literatur (\S12), dan daftar klaim yang harus
di-\emph{tone-down} (\S11.3). Rekomendasi: tandai file tersebut sebagai arsip historis dengan
penunjuk ke dokumen ini.

\noindent
\textbf{Satu koreksi diagnosis.} Catatan kerja sebelumnya mendiagnosis dua tes gagal ini sebagai
``config drift'' pada bobot Ramadan dan import policy. Diagnosis itu \textbf{keliru}. Nilai bobot
di \code{scoring.py} konsisten dan berjumlah 1,00; penyebab sebenarnya adalah ketergantungan
tanggal yang diuraikan di \S\ref{sec:f1}.

% ---------------------------------------------------------------
\subsection{F6 (\sevL) API produksi berjalan dalam mock mode}
\label{sec:f6}

Endpoint \code{/health} melaporkan \code{mock\_mode: true} dan \code{gemini\_mock: true}. Untuk
demo hackathon ini wajar dan aman (tanpa kredensial, tanpa biaya, deterministik). Dicatat semata
agar tidak ada yang menyangka jalur Gemini dan Twilio live sudah tervalidasi di produksi. Yang
tervalidasi hidup adalah mesin pencocokan, data, dan penyajiannya.

% ---------------------------------------------------------------
\subsection{F7 (\sevL) Parameter city pada endpoint forecast tidak intuitif}
\label{sec:f7}

\code{/api/v1/forecast} mensyaratkan \code{city} berupa kode kabupaten BPS, bukan nama kota.
Permintaan \code{city=Surabaya} gagal, sedangkan \code{city=3578} berhasil. Mitigasinya sudah
baik: respons error mengembalikan daftar \code{available\_pairs}, sehingga dapat ditemukan sendiri
oleh pengguna. Perbaikan opsional adalah menerima nama kota sebagai alias.

% =====================================================================
\section{Klaim versus Bukti}
\label{sec:klaim}

Tabel ini menguji setiap klaim menonjol di README terhadap bukti terukur.

\begin{center}
\begin{tabular}{@{}p{5.6cm}cp{7.8cm}@{}}
\toprule
\textbf{Klaim README} & \textbf{Status} & \textbf{Bukti} \\
\midrule
403 tes lulus, 1 di-skip & \ok & Direproduksi persis di lokal \\
Mesin 4 lapis di data BPS asli 2022 & \ok & 63 match, 415.210 ton \\
Deteksi anomali pada PIHPS 2021--2025 & \ok & Berjalan, label metode salah (F3) \\
Forecasting dengan TimesFM 2.0 & \gagal & 56/56 artefak baseline (F2) \\
Dashboard peta interaktif & \ok & HTTP 200, terhubung ke API produksi \\
Bot WhatsApp & \ok & Kode ada, produksi dalam mock mode (F6) \\
Equity terbukti saat pasokan langka & \ok & Fixture constrained ter-commit \\
Biaya efisiensi nol & \ok & Coverage identik 0,6649, Gini 0,3017 ke 0,2905 \\
Tech stack sengaja diringkas & \ok & Konsisten dengan isi repositori \\
\bottomrule
\end{tabular}
\end{center}

\subsection{Catatan tentang klaim equity}

Klaim equity AgriFlow terverifikasi \textbf{dan} dibingkai dengan jujur. README secara eksplisit
menyatakan ``Kami tidak mengklaim keunggulan equity saat pasokan melimpah'', yang memang sesuai
bukti: pada skenario abundant, nilai equity tidak terlihat, dan Gini AgriFlow (0,099) bahkan
sedikit lebih buruk daripada greedy (0,093). Pembingkaian pada skenario supply-constrained adalah
pilihan yang tepat dan dapat dipertahankan.

\subsection{Catatan tentang performa}

\begin{center}
\begin{tabular}{@{}lrrrr@{}}
\toprule
\textbf{Konfigurasi} & \textbf{N (s$\times$d)} & \textbf{p50 (ms)} & \textbf{p99 (ms)} & \textbf{vs target 500 ms} \\
\midrule
Sample data CSV (realistis) & 40$\times$33   & 0,72  & 1,44  & 347$\times$ margin \\
Sintetis full Jatim         & 361$\times$361 & 41,62 & 59,53 & 8,4$\times$ margin \\
Stress 100$\times$100       & 100$\times$100 & 8,86  & 15,32 & 33$\times$ margin \\
Stress 200$\times$200       & 200$\times$200 & 20,90 & 34,09 & 15$\times$ margin \\
\bottomrule
\end{tabular}
\end{center}

\noindent
Seluruh konfigurasi lulus target 500 ms dengan p99 tertinggi 59,53 ms (margin 88,1\,\%). Angka ini
mengonfirmasi ulang temuan \code{AUDIT\_v10.md}: skala provinsi aman, sedangkan skala nasional 514
kabupaten tetap memerlukan spatial indexing sebagaimana didokumentasikan di sana.

% =====================================================================
\section{Hal yang Terverifikasi Baik}

Agar audit ini berimbang, berikut yang terbukti solid dan sebaiknya tidak diutak-atik:

\begin{itemize}
  \item \textbf{Kebersihan kredensial.} Tidak ada \code{.env} ter-track, tidak ada secret di
        seluruh riwayat commit, dan \code{.gitignore} memiliki aturan defensif eksplisit
        (\code{.secrets}, \code{secrets/}, \code{*.key}).
  \item \textbf{CORS dikonfigurasi ketat.} Origin di-whitelist spesifik, bukan wildcard.
  \item \textbf{Penjaga OOM pada jalur forecast.} Server tidak pernah meng-import TimesFM saat
        runtime, dan fallback diberi label jujur di payload API.
  \item \textbf{Penanganan kasus tepi yang anggun.} \code{bawang\_putih} tanpa surplus domestik
        tidak membuat mesin gagal, melainkan menghasilkan saran impor.
  \item \textbf{Matriks CI lintas OS dan lintas versi Python} sudah ada dan berjalan, sehingga
        kegagalan tanggal ini justru tertangkap. Ini pertanda CI bekerja sebagaimana mestinya.
  \item \textbf{Asumsi demo dinyatakan terbuka.} Keluaran \code{run\_demo\_real.py} mencantumkan
        sendiri asumsinya (margin 25\,\%, \code{harvest\_age\_days}), bukan menyembunyikannya.
\end{itemize}

% =====================================================================
\section{Rekomendasi Berprioritas}

\begin{center}
\begin{tabular}{@{}clp{7.4cm}ll@{}}
\toprule
\textbf{\#} & \textbf{Prioritas} & \textbf{Tindakan} & \textbf{Effort} & \textbf{Ref} \\
\midrule
1 & \sevH & Perbaiki prioritas event dan pisahkan kebijakan impor dari event kalender & 2--4 jam & F1 \\
2 & \sevH & Sematkan \code{reference\_date} pada 2 tes agar deterministik & 30 menit & F1 \\
3 & \sevH & Koreksi baris forecasting di README menjadi label baseline & 15 menit & F2 \\
4 & \sevM & Push perbaikan pelabelan detektor anomali ke produksi & 1 jam & F3 \\
5 & \sevM & Cadangkan branch \code{post-submit-analysis} ke remote (perlu izin) & 10 menit & F4 \\
6 & \sevM & Tandai \code{AUDIT\_v10.md} sebagai arsip, tunjuk ke dokumen ini & 15 menit & F5 \\
7 & \sevL & Terima nama kota sebagai alias pada parameter \code{city} & 1 jam & F7 \\
8 & \sevL & Jadikan \code{NEXT\_PUBLIC\_API\_URL} wajib saat build produksi & 30 menit & \S7.1 \\
\bottomrule
\end{tabular}
\end{center}

\noindent
Butir 1 sampai 3 mengembalikan CI ke hijau sekaligus menutup dua overclaim, dengan total usaha di
bawah satu hari kerja.

% =====================================================================
\section{Reproduksi Audit}

Seluruh temuan di dokumen ini dapat direproduksi ulang dengan perintah berikut.

\begin{lstlisting}
# Test suite lengkap
python -m pytest tests/ -q
# Harapan: 403 passed, 1 skipped

# Reproduksi akar masalah F1 (bug kalender)
python -c "from datetime import datetime; \
from matching_engine.engine import get_active_demand_event; \
print(get_active_demand_event(datetime(2026,7,14)))"
# Harapan: SCHOOL_START  (tanggal CI yang gagal)

# Reproduksi F1 pada test suite: jalankan dua tes ini antara 1-15 Jul / 1-15 Jan
python -m pytest tests/test_scenarios_political.py::TestE4_ImportPolicy \
                 tests/test_scenarios_temporal.py::TestC1_RamadanSpike -v

# Reproduksi F2 (artefak forecast berlabel baseline)
python -c "import json,collections; \
d=json.load(open('sample_data/forecasts/forecast_all.json')); \
print(collections.Counter(r['method'] for r in d))"
# Harapan: Counter({'seasonal_naive_baseline': 56})

# Demo data nyata BPS 2022
python examples/run_demo_real.py

# Benchmark latency
python benchmarks/latency.py

# Log CI yang gagal
gh run view 29348620529 --log-failed

# Verifikasi API produksi
curl https://masteraaa123-agriflow-api.hf.space/health
curl "https://masteraaa123-agriflow-api.hf.space/api/v1/forecast?commodity=beras_premium&city=3578"
\end{lstlisting}

% =====================================================================
\section{Kesimpulan}

AgriFlow adalah proyek yang secara teknis nyata: mesin pencocokannya berjalan di data BPS asli,
API-nya hidup di produksi, dashboard-nya tersambung end-to-end, dan 403 tes melindunginya dari
regresi. Fondasi rekayasanya kuat, dan budaya \emph{honest engineering} yang diklaim README
sebagian besar benar-benar tercermin di kode, khususnya pada pelabelan fallback forecast yang
jujur dan pembingkaian klaim equity yang disiplin.

Yang perlu dibereskan bukan arsitekturnya, melainkan tiga titik kecil di lapisan integritas: satu
cacat logika kalender yang membuat CI merah dan berpotensi menghasilkan jejak audit yang
kontradiktif, serta dua klaim dokumentasi yang melampaui artefak yang benar-benar dikirim. Menarik
dicatat bahwa pada kedua kasus overclaim, kode dan API justru sudah jujur; hanya narasi di README
yang perlu diturunkan agar setara dengan bukti.

Ketiganya dapat diselesaikan dalam waktu kurang dari satu hari kerja. Setelah itu, seluruh klaim
proyek ini akan sepenuhnya dapat dipertahankan di hadapan pemeriksaan teknis mana pun.

\vspace{1.5em}
\hrule
\vspace{0.6em}
\noindent\small\textit{Audit dijalankan pada 17 Juli 2026 terhadap commit \code{a6b3b47}. Setiap
temuan diverifikasi melalui eksekusi kode, pemanggilan endpoint produksi, atau pembacaan log CI.
Tidak ada perubahan yang diterapkan ke kode, repositori, maupun deployment selama audit ini.}

\end{document}
